%!TeX root=../persuasiontop.tex
\chapter[Chapter \thechapter]{}

\begin{figure}[t!]
\centering
\includegraphics[width=\linewidth]{nouveau14}
\end{figure}

\lettrine[lraise=0.3]{T}{hough} Charles and Mary had remained at Lyme much longer after Mr and Mrs Musgrove's going than Anne conceived they could have been at all wanted, they were yet the first of the family to be at home again; and as soon as possible after their return to Uppercross they drove over to the Lodge. They had left Louisa beginning to sit up; but her head, though clear, was exceedingly weak, and her nerves susceptible to the highest extreme of tenderness; and though she might be pronounced to be altogether doing very well, it was still impossible to say when she might be able to bear the removal home; and her father and mother, who must return in time to receive their younger children for the Christmas holidays, had hardly a hope of being allowed to bring her with them.

They had been all in lodgings together. Mrs Musgrove had got Mrs Harville's children away as much as she could, every possible supply from Uppercross had been furnished, to lighten the inconvenience to the Harvilles, while the Harvilles had been wanting them to come to dinner every day; and in short, it seemed to have been only a struggle on each side as to which should be most disinterested and hospitable.

Mary had had her evils; but upon the whole, as was evident by her staying so long, she had found more to enjoy than to suffer. Charles Hayter had been at Lyme oftener than suited her; and when they dined with the Harvilles there had been only a maid-servant to wait, and at first Mrs Harville had always given Mrs Musgrove precedence; but then, she had received so very handsome an apology from her on finding out whose daughter she was, and there had been so much going on every day, there had been so many walks between their lodgings and the Harvilles, and she had got books from the library, and changed them so often, that the balance had certainly been much in favour of Lyme. She had been taken to Charmouth too, and she had bathed, and she had gone to church, and there were a great many more people to look at in the church at Lyme than at Uppercross; and all this, joined to the sense of being so very useful, had made really an agreeable fortnight.

Anne enquired after Captain Benwick. Mary's face was clouded directly. Charles laughed.

<Oh! Captain Benwick is very well, I believe, but he is a very odd young man. I do not know what he would be at. We asked him to come home with us for a day or two: Charles undertook to give him some shooting, and he seemed quite delighted, and, for my part, I thought it was all settled; when behold! on Tuesday night, he made a very awkward sort of excuse; <he never shot> and he had <been quite misunderstood,> and he had promised this and he had promised that, and the end of it was, I found, that he did not mean to come. I suppose he was afraid of finding it dull; but upon my word I should have thought we were lively enough at the Cottage for such a heart-broken man as Captain Benwick.>

Charles laughed again and said, <Now Mary, you know very well how it really was. It was all your doing,> (turning to Anne). <He fancied that if he went with us, he should find you close by: he fancied everybody to be living in Uppercross; and when he discovered that Lady Russell lived three miles off, his heart failed him, and he had not courage to come. That is the fact, upon my honour. Mary knows it is.>

But Mary did not give into it very graciously, whether from not considering Captain Benwick entitled by birth and situation to be in love with an Elliot, or from not wanting to believe Anne a greater attraction to Uppercross than herself, must be left to be guessed. Anne's good-will, however, was not to be lessened by what she heard. She boldly acknowledged herself flattered, and continued her enquiries.

<Oh! he talks of you,> cried Charles, <in such terms—> Mary interrupted him. <I declare, Charles, I never heard him mention Anne twice all the time I was there. I declare, Anne, he never talks of you at all.>

<No,> admitted Charles, <I do not know that he ever does, in a general way; but however, it is a very clear thing that he admires you exceedingly. His head is full of some books that he is reading upon your recommendation, and he wants to talk to you about them; he has found out something or other in one of them which he thinks—oh! I cannot pretend to remember it, but it was something very fine—I overheard him telling Henrietta all about it; and then <Miss Elliot> was spoken of in the highest terms! Now Mary, I declare it was so, I heard it myself, and you were in the other room. <Elegance, sweetness, beauty.> Oh! there was no end of Miss Elliot's charms.>

<And I am sure,> cried Mary, warmly, <it was a very little to his credit, if he did. Miss Harville only died last June. Such a heart is very little worth having; is it, Lady Russell? I am sure you will agree with me.>

<I must see Captain Benwick before I decide,> said Lady Russell, smiling.

<And that you are very likely to do very soon, I can tell you, ma'am,> said Charles. <Though he had not nerves for coming away with us, and setting off again afterwards to pay a formal visit here, he will make his way over to Kellynch one day by himself, you may depend on it. I told him the distance and the road, and I told him of the church's being so very well worth seeing; for as he has a taste for those sort of things, I thought that would be a good excuse, and he listened with all his understanding and soul; and I am sure from his manner that you will have him calling here soon. So, I give you notice, Lady Russell.>

<Any acquaintance of Anne's will always be welcome to me,> was Lady Russell's kind answer.

<Oh! as to being Anne's acquaintance,> said Mary, <I think he is rather my acquaintance, for I have been seeing him every day this last fortnight.>

<Well, as your joint acquaintance, then, I shall be very happy to see Captain Benwick.>

<You will not find anything very agreeable in him, I assure you, ma'am. He is one of the dullest young men that ever lived. He has walked with me, sometimes, from one end of the sands to the other, without saying a word. He is not at all a well-bred young man. I am sure you will not like him.>

<There we differ, Mary,> said Anne. <I think Lady Russell would like him. I think she would be so much pleased with his mind, that she would very soon see no deficiency in his manner.>

<So do I, Anne,> said Charles. <I am sure Lady Russell would like him. He is just Lady Russell's sort. Give him a book, and he will read all day long.>

<Yes, that he will!> exclaimed Mary, tauntingly. <He will sit poring over his book, and not know when a person speaks to him, or when one drops one's scissors, or anything that happens. Do you think Lady Russell would like that?>

Lady Russell could not help laughing. <Upon my word,> said she, <I should not have supposed that my opinion of any one could have admitted of such difference of conjecture, steady and matter of fact as I may call myself. I have really a curiosity to see the person who can give occasion to such directly opposite notions. I wish he may be induced to call here. And when he does, Mary, you may depend upon hearing my opinion; but I am determined not to judge him beforehand.>

<You will not like him, I will answer for it.>

Lady Russell began talking of something else. Mary spoke with animation of their meeting with, or rather missing, Mr Elliot so extraordinarily.

<He is a man,> said Lady Russell, <whom I have no wish to see. His declining to be on cordial terms with the head of his family, has left a very strong impression in his disfavour with me.>

This decision checked Mary's eagerness, and stopped her short in the midst of the Elliot countenance.

With regard to Captain Wentworth, though Anne hazarded no enquiries, there was voluntary communication sufficient. His spirits had been greatly recovering lately as might be expected. As Louisa improved, he had improved, and he was now quite a different creature from what he had been the first week. He had not seen Louisa; and was so extremely fearful of any ill consequence to her from an interview, that he did not press for it at all; and, on the contrary, seemed to have a plan of going away for a week or ten days, till her head was stronger. He had talked of going down to Plymouth for a week, and wanted to persuade Captain Benwick to go with him; but, as Charles maintained to the last, Captain Benwick seemed much more disposed to ride over to Kellynch.

There can be no doubt that Lady Russell and Anne were both occasionally thinking of Captain Benwick, from this time. Lady Russell could not hear the door-bell without feeling that it might be his herald; nor could Anne return from any stroll of solitary indulgence in her father's grounds, or any visit of charity in the village, without wondering whether she might see him or hear of him. Captain Benwick came not, however. He was either less disposed for it than Charles had imagined, or he was too shy; and after giving him a week's indulgence, Lady Russell determined him to be unworthy of the interest which he had been beginning to excite.

The Musgroves came back to receive their happy boys and girls from school, bringing with them Mrs Harville's little children, to improve the noise of Uppercross, and lessen that of Lyme. Henrietta remained with Louisa; but all the rest of the family were again in their usual quarters.

Lady Russell and Anne paid their compliments to them once, when Anne could not but feel that Uppercross was already quite alive again. Though neither Henrietta, nor Louisa, nor Charles Hayter, nor Captain Wentworth were there, the room presented as strong a contrast as could be wished to the last state she had seen it in.

Immediately surrounding Mrs Musgrove were the little Harvilles, whom she was sedulously guarding from the tyranny of the two children from the Cottage, expressly arrived to amuse them. On one side was a table occupied by some chattering girls, cutting up silk and gold paper; and on the other were tressels and trays, bending under the weight of brawn and cold pies, where riotous boys were holding high revel; the whole completed by a roaring Christmas fire, which seemed determined to be heard, in spite of all the noise of the others. Charles and Mary also came in, of course, during their visit, and Mr Musgrove made a point of paying his respects to Lady Russell, and sat down close to her for ten minutes, talking with a very raised voice, but from the clamour of the children on his knees, generally in vain. It was a fine family-piece.

Anne, judging from her own temperament, would have deemed such a domestic hurricane a bad restorative of the nerves, which Louisa's illness must have so greatly shaken. But Mrs Musgrove, who got Anne near her on purpose to thank her most cordially, again and again, for all her attentions to them, concluded a short recapitulation of what she had suffered herself by observing, with a happy glance round the room, that after all she had gone through, nothing was so likely to do her good as a little quiet cheerfulness at home.

Louisa was now recovering apace. Her mother could even think of her being able to join their party at home, before her brothers and sisters went to school again. The Harvilles had promised to come with her and stay at Uppercross, whenever she returned. Captain Wentworth was gone, for the present, to see his brother in Shropshire.

<I hope I shall remember, in future,> said Lady Russell, as soon as they were reseated in the carriage, <not to call at Uppercross in the Christmas holidays.>

Everybody has their taste in noises as well as in other matters; and sounds are quite innoxious, or most distressing, by their sort rather than their quantity. When Lady Russell not long afterwards, was entering Bath on a wet afternoon, and driving through the long course of streets from the Old Bridge to Camden Place, amidst the dash of other carriages, the heavy rumble of carts and drays, the bawling of newspapermen, muffin-men and milkmen, and the ceaseless clink of pattens, she made no complaint. No, these were noises which belonged to the winter pleasures; her spirits rose under their influence; and like Mrs Musgrove, she was feeling, though not saying, that after being long in the country, nothing could be so good for her as a little quiet cheerfulness.

Anne did not share these feelings. She persisted in a very determined, though very silent disinclination for Bath; caught the first dim view of the extensive buildings, smoking in rain, without any wish of seeing them better; felt their progress through the streets to be, however disagreeable, yet too rapid; for who would be glad to see her when she arrived? And looked back, with fond regret, to the bustles of Uppercross and the seclusion of Kellynch.

Elizabeth's last letter had communicated a piece of news of some interest. Mr Elliot was in Bath. He had called in Camden Place; had called a second time, a third; had been pointedly attentive. If Elizabeth and her father did not deceive themselves, had been taking much pains to seek the acquaintance, and proclaim the value of the connection, as he had formerly taken pains to shew neglect. This was very wonderful if it were true; and Lady Russell was in a state of very agreeable curiosity and perplexity about Mr Elliot, already recanting the sentiment she had so lately expressed to Mary, of his being <a man whom she had no wish to see.> She had a great wish to see him. If he really sought to reconcile himself like a dutiful branch, he must be forgiven for having dismembered himself from the paternal tree.

Anne was not animated to an equal pitch by the circumstance, but she felt that she would rather see Mr Elliot again than not, which was more than she could say for many other persons in Bath.

She was put down in Camden Place; and Lady Russell then drove to her own lodgings, in Rivers Street.